\section{Introduction}
\textbf{Motivation.} Dịch vụ dựa trên vị trí (LBS) giúp người dùng tìm địa điểm quan tâm (POI) khi di chuyển.
Những yêu cầu liên tiếp đồng thời tạo thành dấu vết mà nhà cung cấp có thể kết hợp với bản đồ và lịch sử
để suy luận vị trí, đường đi hoặc các thông tin liên quan. Bảo vệ một điểm gửi riêng lẻ vì vậy chưa đủ
để mô tả rủi ro của cả quá trình sử dụng dịch vụ. Các nghiên cứu về bối cảnh giao thông, ngữ nghĩa và
query liên tục đã khai thác những tương quan này\cite{transprotect,semantic,fake}.

\textbf{Các vấn đề cần giải quyết.} Thứ nhất, cơ chế nhiều bước phải tính cả quyết định khi nào cập nhật
thông tin vị trí, không chỉ việc thêm nhiễu. Thứ hai, đầu ra giả có thể hợp lý về hình học nhưng không
lấy được POI phù hợp cho người dùng. Thứ ba, các công trình được khảo sát dùng output, attacker, dữ liệu
và metrics khác nhau: một tọa độ thay thế, một tập chứa GPS thật và các dummy, hoặc một bộ dữ liệu
trajectory được tổng hợp offline không thể được so sánh bằng một score không đổi ý nghĩa.

Vấn đề đánh giá không nằm ở việc hoàn toàn thiếu dataset hoặc metrics trong lĩnh vực. Điều còn thiếu
cho phép đối chiếu ở đây là \emph{một protocol thống nhất giữa những hợp đồng đầu ra và mục tiêu đang xét}.
Luận văn do đó hệ thống hóa các tác vụ suy luận, xây dựng workload có kiểm soát và sử dụng hai lớp đo:
metrics chung theo tác vụ, cùng metrics gốc khi định nghĩa và đầu vào tương thích.

\begin{tabularx}{\linewidth}{@{}P{2.1cm}P{5.6cm}Y@{}}\toprule
\textbf{Nhóm} & \textbf{Mục tiêu} & \textbf{Vai trò nghiên cứu}\\\midrule
S1--S3 & Vị trí hiện tại, nơi dừng, đường đã đi & Tác vụ định vị/tái dựng theo các phạm vi quan sát.\\
S4--S6 & Liên kết người/xe, cạnh kế tiếp, đích tương lai & Suy luận phụ thuộc prior và quan hệ giữa target với trajectory.\\
S7--S8 & Nội dung/ý định query; thông tin người đi cùng & Phân biệt payload với intent, cá nhân với tương quan nhóm.\\
S9--S10 & Điểm đầu và điểm cuối phiên & Rủi ro endpoint dưới mô hình quan sát được khai báo.\\\bottomrule
\end{tabularx}

\textbf{Hướng giải quyết.} GeoI-BQ giữ Geo-indistinguishability (Geo-I) làm nền tảng. Thiết bị dùng REM
và phép thử tái sử dụng có nhiễu để duy trì điểm tham chiếu Z. Từ lịch sử đã bảo vệ, một belief gần đúng
đánh giá những vị trí người dùng có thể ở. Belief điều chỉnh phân phối chọn \emph{đoạn Q} trong giới hạn;
các khung Q được phát theo lịch, còn GPS hiện tại và nhu cầu thật chỉ dùng tại bộ lọc/xếp hạng local.
Phân tích toán học nối các thao tác này thành cận cho toàn transcript trước khi đọc benchmark.

\textbf{Đóng góp chính.}
\begin{enumerate}
\item Hệ thống hóa bài toán thành các scenario, ranh giới quan sát và protocol đánh giá theo output contract;
xây dựng workload SUMO/POI để đo suy luận vị trí, chất lượng đáp án và chi phí trên cùng dữ liệu.
\item Đặc tả phương pháp dựa trên Geo-I và cơ chế chọn chung đoạn Q; phân tích privacy toàn transcript,
giới hạn suy luận và các bảo đảm utility có điều kiện. Không gọi các nền tảng hợp thành/hậu xử lý là định lý mới.
\item Đối chiếu GeoI-BQ với ba adaptation từ nghiên cứu gần đây trên cùng cohort, kèm Raw và REM-Z controls;
giữ CI, failure, metrics không áp dụng và sự khác biệt về chi phí/bảo đảm.
\end{enumerate}
Các kết quả không mặc định bảo vệ đầy đủ S1--S10. Bằng chứng thực nghiệm chính là S1--S3 và S9--S10;
S4--S6 nhận phân tích suy luận có điều kiện, S7 có tính bất biến payload và S8 là hướng mở rộng.

Phần 2 khảo sát nghiên cứu liên quan; Phần 3 trình bày hệ thống, thuật toán và chứng minh; Phần 4 định nghĩa
đánh giá và phân tích thực nghiệm; Phần 5 tổng kết đóng góp, giới hạn và hướng tiếp theo.
